% Generated draft: requires mathematical and source review.
\documentclass[11pt]{article}
\usepackage[margin=25mm]{geometry}
\usepackage{fontspec}
\setmainfont{Latin Modern Roman}
\usepackage{amsmath,amssymb,mathtools,unicode-math}
\setmathfont{Latin Modern Math}
\usepackage{enumitem,needspace,xurl,hyperref}
\setlength{\parindent}{0pt}
\setlength{\parskip}{5pt}
\setlength{\emergencystretch}{3em}
\newcommand{\N}{\mathbb{N}}
\newcommand{\Z}{\mathbb{Z}}
\newcommand{\Q}{\mathbb{Q}}
\newcommand{\R}{\mathbb{R}}
\newcommand{\C}{\mathbb{C}}
\newcommand{\F}{\mathbb{F}}
\newcommand{\A}{\mathbb{A}}
\newcommand{\PP}{\mathbb{P}}
\newcommand{\E}{\mathbb{E}}
\newcommand{\eps}{\varepsilon}
\newcommand{\dd}{\,\mathrm{d}}
\begin{document}
\section*{UnsolvedMath ID: 30006230}
\textbf{Source problem code:} catalogue entry
\subsection{Definitions and mathematical statement}
Let \((V,W)\) denote the pair of weights appearing in the supplied Grushin-operator problem. For an intended Grushin setting \(G\), write \(\mathsf{H}_{G}(V,W)\) for validity of the associated Hardy inequality, \(\mathsf{HR}_{G}(V,W)\) for validity of the associated Hardy--Rellich inequality, and \(\mathsf{B}_{G}(V,W)\) for the associated Bessel-pair condition. These are placeholders for source-specific conditions, not definitions of those conditions: the supplied record gives neither the inequalities nor the differential equation and positivity requirement defining the Bessel pair.

The requested converse question is whether, for every admissible pair of weights in the intended setting, the relevant implication holds:
\[
\mathsf{H}_{G}(V,W)\Longrightarrow\mathsf{B}_{G}(V,W),
\qquad\text{or, for the intended Hardy--Rellich version,}\qquad
\mathsf{HR}_{G}(V,W)\Longrightarrow\mathsf{B}_{G}(V,W).
\]
Each intended converse is to be proved or disproved under the hypotheses of its corresponding source inequality. The record does not determine whether these are separate questions or which precise variants are intended. In particular, it does not specify the Grushin vector fields, dimensions, operator parameters, domain, measure, admissibility and regularity of the weights, test-function class, boundary or support conditions, constants, or Bessel-pair equation. Consequently, the displayed implications record the logical target only; a fully specified mathematical statement requires those missing source hypotheses.
\subsection{Short English statement}
Does validity of the relevant Hardy or Hardy--Rellich inequality for Grushin vector fields force its weight pair \((V,W)\) to satisfy the corresponding Bessel-pair condition? The supplied excerpt leaves the exact inequalities and hypotheses unspecified, so the converse remains a statement template requiring human mathematical review.
\subsection{Sources}
\begin{itemize}
\item Local UnsolvedMath registry, permanent ID \(30006230\), title: \emph{Converse Bessel-Pair Criterion for Grushin Operators}; exact source code: \texttt{catalogue entry}. The supplied registry labels the problem open and its statement \texttt{corrected\_verified}, and reports that reconstruction was checked against the official report. These are supplied metadata, not independent verification here.
\item \emph{Mini-Workshop: Hardy Inequalities in Discrete and Continuum Settings} (2025), Oberwolfach Reports, DOI 10.4171/owr/2025/14, pp.\ 579--614: \url{https://doi.org/10.4171/owr/2025/14}. The supplied original excerpt identifies p.\ 607 and the contribution \emph{Hardy and Rellich identities for Grushin operators}, presented by Prasun Roychowdhury, joint work with Debdip Ganguly and Jotsaroop Kaur. The registry supplies proposed year 2025 but no separate proposer.
\item Supplied literature reference: Ganguly--Kaur--Roychowdhury, \emph{Grushin Bessel pairs}: \url{https://arxiv.org/abs/2404.05510}.
\item Supplied literature reference: Cossetti--D'Arca, \emph{unified Bessel pairs}: \url{https://arxiv.org/abs/2511.10039}.
\item Literature assessment dated exactly \texttt{2026-08-22}: the supplied curation labels the problem open, states that the Grushin work derives Hardy and Rellich inequalities from a Bessel-pair condition while leaving its converse open, and states that the 2025 general framework does not resolve this converse. This is supplied dated curation derived from the Oberwolfach material, not independently verified; its time-sensitive status and cited claims require independent checking.
\end{itemize}
\end{document}
